\section*{الفصل \ref{ch:b3:point-groups} --- التناظر والزمر النقطية}

\begin{solution}{exo:b3:point-groups:1}
الإيثين: ثلاثة محاور $C_2$ متعامدة، وثلاثة مستويات، و$i$: $D_{2h}$. \ce{XeF4} (مربع
مستوٍ): $C_4$، وأربعة محاور $C_2 \perp C_4$، و$\sigma_h$، و$2\sigma_v$، و$2\sigma_d$، و$i$، و$S_4$:
$D_{4h}$. \ce{PCl5} (هرم مزدوج مثلثي): $C_3$، وثلاثة محاور $C_2$، و$\sigma_h$، و$3\sigma_v$،
و$S_3$: $D_{3h}$. 1،3،5-ثلاثي كلوروبنزين: العناصر نفسها التي للجزيء \ce{BF3}: $D_{3h}$.
\end{solution}

\begin{solution}{exo:b3:point-groups:2}
الإيثان المتخالف: $C_3$، وثلاثة محاور $C_2$ عمودية، وثلاثة مستويات $\sigma_d$، و$i$، و$S_6$:
$D_{3d}$. المتطابق: $C_3$، وثلاثة محاور $C_2$، و$\sigma_h$، وثلاثة مستويات $\sigma_v$: $D_{3h}$.
الهكسان الحلقي في تشكّل الكرسي: $D_{3d}$.
\end{solution}

\begin{solution}{exo:b3:point-groups:3}
\ce{SO3} ($D_{3h}$): لا. \ce{SO2} ($C_{2v}$): نعم. \ce{PCl3} ($C_{3v}$): نعم.
\ce{XeF4} ($D_{4h}$): لا. \ce{CH2Cl2} ($C_{2v}$): نعم. 1،2-ثنائي كلوروإيثين cis
($C_{2v}$): نعم؛ وtrans ($C_{2h}$): لا.
\end{solution}

\begin{solution}{exo:b3:point-groups:4}
$C_2(z) = \mathrm{diag}(-1,-1,1)$، $\chi = -1$؛ $i = \mathrm{diag}(-1,-1,-1)$،
$\chi = -3$؛ $\sigma_h(xy) = \mathrm{diag}(1,1,-1)$، $\chi = 1$.
\end{solution}

\begin{solution}{exo:b3:point-groups:5}
بالمصفوفات القطرية في التمرين السابق: $C_2i = \sigma_h$، $C_2\sigma_h =
i$، $i\sigma_h = C_2$، وكل عملية مقلوب نفسها، وكل جداء
تبديلي. فالزمرة تبديلية، ومنه $X^{-1}AX = A$ من أجل كل $X$: فكل
عملية صنف بمفردها (أربعة أصناف، وأربعة \omterm{def:b3:point-groups:irrep}{تمثيلات غير قابلة للاختزال}
بعدها 1).
\end{solution}

\begin{solution}{exo:b3:point-groups:6}
لنأخذ $\sigma_v(1)$ المستوي $xz$. تنتقل النقطة $(1,0)$ بالدوران $C_3$ إلى $(-\frac12,
\frac{\sqrt3}2)$، ثم بالانعكاس $\sigma_v(1)$ إلى $(-\frac12,-\frac{\sqrt3}2)$؛ وبالترتيب
الآخر تنتقل إلى $(1,0)$ ثم إلى $(-\frac12,\frac{\sqrt3}2)$. فالجداءان
انعكاسان في مستويين مختلفين ($\sigma_v(3)$ و$\sigma_v(2)$): الزمرة
غير تبديلية، ولهذا تشكّل انعكاساتها صنفًا من ثلاث عمليات.
\end{solution}

\begin{solution}{exo:b3:point-groups:7}
يحفظ ثنائي الفينيل الملتوي المحور $C_2$ على امتداد الرابطة بين الحلقتين ومحورين $C_2$
عموديين عليه، ينصّفان الزاويتين بين مستويي الحلقتين؛ وتضيع كل
المستويات \omterm{def:b3:point-groups:inversion}{ومركز الانقلاب} التي للصورتين المستوية أو المتعامدة.
فبثلاثة محاور $C_2$ متعامدة ودون أي عنصر غير فعلي تكون الزمرة
$D_2$: والجزيء كيرالي (صورتاه الملتويتان متماكبان مرآويان،
قابلان للفصل حين تمنع مجموعات استبدال ضخمة في موضع ortho الدوران).
\end{solution}

\begin{solution}{exo:b3:point-groups:8}
$d_{z^2}$ و$d_{x^2-y^2}$: $a_1$ ($x^2$ و$y^2$ و$z^2$ من النوع $A_1$)؛ $d_{xy}$: $a_2$؛ $d_{xz}$:
$b_1$؛ $d_{yz}$: $b_2$.
\end{solution}

\begin{solution}{exo:b3:point-groups:9}
$1 + 1 + 4 + 9 + 9 = 24 = h$. $\sum g\,\chi_{A_1}\chi_{T_2} = 1\cdot3 + 8\cdot0 + 3(-1)
+ 6(-1) + 6(1) = 0$.
\end{solution}

\begin{solution}{exo:b3:point-groups:10}
في المستوي العمودي على خط التقاطع، ينقل الانعكاس في مستقيم
زاويته $\alpha$ الزاوية القطبية $\phi$ إلى $2\alpha - \phi$. وانعكاسان، عند
$\alpha$ ثم $\alpha + \theta$: $\phi \to 2\alpha - \phi \to 2(\alpha + \theta) - (2\alpha - \phi) =
\phi + 2\theta$، أي دوران بزاوية $2\theta$ حول الخط. فالمستويان الرأسيان اللذان بينهما
$60^\circ$ يولّدان دورانًا بزاوية $120^\circ$: أي $C_3$.
\end{solution}

\begin{solution}{exo:b3:point-groups:11}
المحور C=C=C محور $C_2$ ومحور $S_4$ (دوران بزاوية $90^\circ$، يبادل
مستويي \ce{CH2}، ثم انعكاس في المستوي المار بالكربون المركزي)؛ 
ومستويا \ce{CH2} من النوع $\sigma_d$؛ ومحوران $C_2$ عموديان على المحور ينصّفانهما.
الرتبة 8: $D_{2d}$، لاكيرالي. وفي \ce{ClHC=C=CHCl} يُزيل الاستبدال المستويات و$S_4$ 
والمحور $C_2$ المحوري؛ ولا يبقى إلا محور $C_2$ واحد عمودي على
المحور: $C_2$، كيرالي --- وهو ألين ذو كيرالية محورية.
\end{solution}

\begin{solution}{exo:b3:point-groups:12}
\ce{SF6}: $O_h$ ($h = 48$). \ce{SF5Cl}: المحور $C_4$ المار بالذرة Cl وأربعة مستويات $\sigma_v$:
$C_{4v}$ ($h = 8$). trans-\ce{SF4Cl2}: $D_{4h}$ ($h = 16$). cis-\ce{SF4Cl2}: $C_{2v}$ ($h
= 4$). كل منها يحفظ مجموعة جزئية من عمليات $O_h$، مغلقة بالنسبة إلى الجداء،
و8 و16 و4 تقسم 48. القطبية: \ce{SF5Cl} وcis-\ce{SF4Cl2}.
\end{solution}

\begin{solution}{pb:b3:point-groups:1}
\textbf{1.} على امتداد الأقطار الأربعة الكبرى المارة بالذرة C وبكل ذرة H؛ يعطي كل منها
$C_3$ و$C_3^2$: 8 دورانات.
\textbf{2.} على امتداد $x$ و$y$ و$z$ عبر مراكز الأوجه: ينقل $C_2(z)$ النقطة $(1,1,1)
\to (-1,-1,1)$، وهي ذرة هيدروجين. والدوران بزاوية $90^\circ$ حول $z$ متبوعًا
بالانعكاس في $xy$ ينقل $(1,1,1) \to (-1,1,1) \to (-1,1,-1)$، وهي ذرة هيدروجين: أي
$S_4$؛ ومع $S_4^3$، 6 عمليات.
\textbf{3.} المستويات $x = \pm y$ و$y = \pm z$ و$x = \pm z$؛ والمستوي $x = y$ يحتوي
$(1,1,1)$ و$(-1,-1,1)$: فكل مستوٍ يحتوي رابطتين C--H.
\textbf{4.} $1 + 8 + 3 + 6 + 6 = 24$.
\textbf{5.} ينقل الانقلاب $(1,1,1)$ إلى $(-1,-1,-1)$، وهي ليست موضع ذرة هيدروجين:
فلا $i$.
\textbf{6.} $E$؛ $8C_3$؛ $3C_2$؛ $6S_4$؛ $6\sigma_d$.
\textbf{7.} \ce{CH3Cl}: $C_{3v}$، $h = 6$.
\textbf{8.} \ce{CH2Cl2}: $C_{2v}$، $h = 4$.
\textbf{9.} \ce{CHCl3}: $C_{3v}$؛ \ce{CCl4}: $T_d$، $h = 24$.
\textbf{10.} \ce{CH2FCl}: $C_s$ ($h = 2$)؛ \ce{CHFClBr}: $C_1$ ($h = 1$).
\textbf{11.} كل منها يحفظ عمليات $T_d$ التي تحترم
الاستبدال؛ و6 و4 و6 و24 و2 و1 كلها تقسم 24.
\textbf{12.} ذرات الهيدروجين الثلاث، تتبادلها $C_3$ و$C_3^2$ والمستويات الثلاثة
$\sigma_v$.
\textbf{13.} كلها عدا \ce{CCl4}: على امتداد المحور $C_3$ للجزيئين \ce{CH3Cl} و\ce{CHCl3}،
والمحور $C_2$ للجزيء \ce{CH2Cl2}، وفي \omterm{def:b3:point-groups:mirror}{مستوي المرآة} للجزيء \ce{CH2FCl}، ودون
اتجاه مفروض للجزيء \ce{CHFClBr}.
\textbf{14.} \ce{CHFClBr} وحده ($C_1$، لا $S_n$).
\textbf{15.} متماكبان مرآويان.
\textbf{16.} ذرتا الهيدروجين فيه متكافئتان: فالمستوي المار بالذرات C وF وCl
والمنصّف للزاوية H--C--H \omterm{def:b3:point-groups:mirror}{مستوي مرآة}.
\textbf{17.} لا: أربع مجموعات استبدال مختلفة لا تُبقي إلا $E$؛ ودون $S_n$ يكون
الجزيء كيراليًا (ويمكن أن يكون قطبيًا).
\textbf{18.} $(x, y, z)$: $T_2$.
\textbf{19.} ذرتا هيدروجين هما طرفا حرف من رباعي الأوجه؛ والعمليات الأربع والعشرون
تنقل أي حرف إلى أي حرف آخر (الأحرف الستة تشكّل مجموعة واحدة): فثمة متماكب
واحد للجزيء \ce{CH2Cl2}.
\textbf{20.} \ce{CH2FCl}: واحد. \ce{CHFClBr}: اثنان، زوج من المتماكبات المرآوية.
\textbf{21.} ثلاثة: ortho ($C_{2v}$)، وmeta ($C_{2v}$)، وpara ($D_{2h}$).
\textbf{22.} ثلاثة: 1،2،3- ($C_{2v}$)، و1،2،4- ($C_s$)، و1،3،5- ($D_{3h}$).
\textbf{23.} اثنان (ذرتا Cl متجاورتان أو متقابلتان). ولا يُعرف إلا \ce{CH2Cl2} واحد،
وهذا استبعد الكربون المستوي ودعم الكربون الرباعي الأوجه
الذي اقتُرح سنة 1874.
\textbf{24.} $1^2 + 1^2 + 2^2 + 3^2 + 3^2 = 24$: \textbf{رتبة $T_d$، \omterm{def:b3:point-groups:point-group}{الزمرة
النقطية} للميثان، هي $h = 24$}.
\end{solution}
